
@@M0@@
il s'%
agit de montrer que les Modules @@M1@@ et @@M2@@
sont canoniquement isomorphes. Or on a d'une part les isomorphismes canoniques :
@@M3@@
D'autre part :
@@M4@@
Enfin, il résulte de EGA II 2.2.2, traduit en termes de faisceaux, qu'on a les isomorphismes canoniques :
@@M5@@
d'où finalement l'isomorphisme cherché.

Il est clair que si @@M6@@ est quasi-cohérent le Module @@M7@@ ainsi construit est bien celui qui est défini dans EGA II.

2.3.3. Il résulte de 2.3.1 et 2.3.2 que le foncteur @@M8@@ est exact. Revenant aux notations de 2.2, on en déduit donc par passage aux catégories dérivées un foncteur encore noté @@M9@@ :
@@M10@@

On voit également sur les constructions précédentes que le foncteur
